%
% Chapter 9 — Conclusion
%
\chapter{Conclusion} \label{cap:conclusion}

\section{Accomplished Work} \label{sec:conc_accomplished}

This beta report documents the progress achieved in the development of the Diário
LX platform. The following components have been implemented:

\begin{itemize}
  \item \textbf{Architecture and data model}: the system architecture was defined,
        the technology stack selected, and the database schema designed and implemented,
        including support for a block-based content model, media credits, category
        hierarchy, and session management.

  \item \textbf{Authentication}: a JWT-based authentication mechanism was implemented,
        with short-lived access tokens and long-lived refresh tokens stored as
        \texttt{HttpOnly} cookies. Token refresh is handled transparently on the
        frontend.

  \item \textbf{Backend API}: the REST API is implemented and covers user management,
        invitation-based onboarding, category and tag management, media upload and
        management, and content creation. The editorial workflow (submission, approval,
        and rejection) is currently under active development.

  \item \textbf{Content editor}: the backoffice content editor is functional,
        supporting block-based authoring with text and image blocks, inline formatting,
        media gallery integration, and metadata configuration (category, tags, authors,
        slug). Support for video content is implemented; podcast and episode creation
        are in an advanced stage of development.

  \item \textbf{Backoffice}: the backoffice interface is largely implemented, covering
        content listing, user management, category and tag management, invitation
        management, and user profiles.

  \item \textbf{Public website}: the article page is implemented and functional.
        The remaining public pages (homepage, category and tag listings, podcast page)
        are under development.

  \item \textbf{Testing and CI}: backend and frontend test suites are in place and
        integrated into a continuous integration pipeline via GitHub Actions.
\end{itemize}

\section{Challenges} \label{sec:conc_challenges}

An early design decision that shaped the data model was the generalisation of the
content entity. The initial implementation was centred around a single
\texttt{article} type; during development, it became clear that the same structure
could naturally accommodate other content types --- videos, podcasts, and episodes
--- with minimal changes. The model was refactored to a generic \texttt{content}
table with a \texttt{type} discriminator, a change that proved straightforward to
implement and significantly improved the long-term scalability of the system.

A second challenge was the selection and integration of the multimedia storage
solution. Evaluating the available options, deploying SeaweedFS, and integrating
it with the backend and frontend required more effort than initially anticipated.
The deployment of the full system to production infrastructure remains a challenge
to be addressed in the final phase.

A third challenge, currently ongoing, is the implementation of the editorial
workflow. Although the state machine itself is well defined, enforcing the full
set of transitions and permission rules consistently across both the backend and
the backoffice interface has revealed additional complexity as development
progresses. The integration between role-based access control and content state
management requires careful coordination, and the approach is being refined
iteratively as the implementation evolves.

The block-based content model, while initially appearing complex, proved to be a
deliberate and effective design choice: by treating the rich text editor (TipTap)
as a tool for editing individual text blocks rather than the entire article body,
the system gains finer control over content structure and media embedding.

\section{Future Work} \label{sec:conc_future}

The following items are planned for the final version:

\begin{itemize}
  \item \textbf{Editorial workflow}: complete the implementation of the
        submission, approval, and rejection workflow in both the backend and
        the backoffice interface, including a dedicated pending content view
        for Editors.

  \item \textbf{Public website}: complete the homepage, category and tag listing
        pages, and the podcast page.

  \item \textbf{Full-text search}: implement keyword search across published
        content, accessible from the public website.

  \item \textbf{Responsive design}: complete the responsive layout across all
        viewports for both the public website and the backoffice.

  \item \textbf{Test coverage}: expand automated test coverage across backend
        and frontend layers.

  \item \textbf{Deployment}: deploy the application to the IPL server
        infrastructure, with a dedicated section documenting the deployment
        process and maintenance procedures.

  \item \textbf{Optional features}: reader comments, scheduled publication,
        and advanced search filters, subject to available time.
\end{itemize}